% =============================================================
% environments.tex — 全 tcolorbox 環境
% =============================================================

% ── 共通バー幅 ──
\newcommand{\BarW}{3.6pt}
\newcommand{\ArcR}{1.5pt}

% ── 汎用ラベル付き左バーボックス ──
\NewDocumentCommand{\MakeActorBox}{m m m m m}{%
  % #1=env name, #2=default label, #3=bar color, #4=bg color, #5=title bg
  \newtcolorbox{#1}[1][#2]{%
    enhanced, breakable,
    colback=#4, colframe=#4, boxrule=0pt,
    leftrule=\BarW,
    colbacktitle=#5, coltitle=white,
    fonttitle=\sffamily\footnotesize\bfseries,
    title=\,##1\,,
    attach boxed title to top left={xshift=-\BarW, yshift=-6pt},
    boxed title style={colframe=#5, boxrule=0pt, sharp corners, arc=0pt,
      left=6pt, right=6pt, top=2pt, bottom=2pt},
    overlay unbroken and first={
      \draw[#3, line width=\BarW] (frame.north west) -- (frame.south west);
    },
    arc=\ArcR, outer arc=0pt, sharp corners=northwest,
    left=12pt, right=10pt, top=14pt, bottom=8pt,
    fontupper=\color{Ink}\small,
    before skip=10pt, after skip=10pt,
  }%
}

% ── 定義/要点 ──
\MakeActorBox{teigi}{要点}{TeigiBar}{TeigiBG}{TeigiBar}

% ── つまずき ──
\MakeActorBox{pitfall}{つまずき}{PitBar}{PitBG}{PitBar}

% ── 比較 ──
\MakeActorBox{hikaku}{比較}{HikakuBar}{HikakuBG}{HikakuBar}

% ── 見抜き方 ──
\MakeActorBox{minuki}{見抜き方}{MinukiBar}{MinukiBG}{MinukiBar}

% ── よくあるミス ──
\MakeActorBox{yokumiss}{よくあるミス}{MissBar}{MissBG}{MissBar}

% ── 解説 ──
\MakeActorBox{kaisetsubox}{解答・解説}{KaiBar}{KaiBG}{KaiBar}

% ── アドバイス ──
\MakeActorBox{advice}{学習アドバイス}{Muted}{LeadBG}{Muted}

% ── 本試験パターン ──
\MakeActorBox{honshiken}{本試験}{MinukiBar}{MinukiBG}{MinukiBar}

% =============================================================
% 例題 (reidai): 紺ヘッダーバー＋問題
% =============================================================
\newcounter{reidaictr}
\renewcommand{\thereidaictr}{\arabic{reidaictr}}

\NewDocumentEnvironment{reidai}{m}{%
  \refstepcounter{reidaictr}%
  \par\addvspace{14pt}%
  \begin{tcolorbox}[
    enhanced, breakable,
    colback=ReidaiBG, colframe=ReidaiBar,
    boxrule=0pt, toprule=3pt,
    arc=0pt, outer arc=0pt,
    left=12pt, right=12pt, top=10pt, bottom=8pt,
    before skip=10pt, after skip=4pt,
    title={\sffamily\small\bfseries 例題 \thereidaictr\quad #1},
    colbacktitle=ReidaiBar, coltitle=white,
    attach boxed title to top left={yshift=0pt, xshift=0pt},
    boxed title style={colframe=ReidaiBar, boxrule=0pt, arc=0pt,
      left=10pt, right=10pt, top=4pt, bottom=4pt},
  ]
  \color{Ink}\small
}{%
  \end{tcolorbox}
  \par\addvspace{4pt}%
}

% ── 指針 ──
\NewDocumentEnvironment{shishin}{}{%
  \begin{tcolorbox}[
    enhanced, breakable,
    colback=ShishinBG, colframe=ShishinBG, boxrule=0pt,
    leftrule=\BarW,
    overlay unbroken and first={
      \draw[ShishinBar, line width=\BarW] (frame.north west) -- (frame.south west);
    },
    arc=\ArcR, outer arc=0pt,
    left=12pt, right=10pt, top=10pt, bottom=8pt,
    before skip=4pt, after skip=4pt,
    fontupper=\color{Ink}\small,
    title={\sffamily\footnotesize\bfseries\color{ShishinBar} 指針},
    coltitle=ShishinBar, colbacktitle=ShishinBG,
    attach boxed title to top left={xshift=8pt, yshift=-2pt},
    boxed title style={colframe=ShishinBG, boxrule=0pt, arc=0pt,
      left=0pt, right=0pt, top=0pt, bottom=0pt},
  ]
}{%
  \end{tcolorbox}
}

% ── 解答（緑左バー）──
\NewDocumentEnvironment{kaitou}{}{%
  \par\addvspace{4pt}%
  \begin{tcolorbox}[
    enhanced, breakable,
    colback=KaitouBG, colframe=KaitouBG, boxrule=0pt,
    leftrule=\BarW,
    colbacktitle=KaitouBar, coltitle=white,
    fonttitle=\sffamily\footnotesize\bfseries,
    title=\,解答\,,
    attach boxed title to top left={xshift=-\BarW, yshift=-6pt},
    boxed title style={colframe=KaitouBar, boxrule=0pt, sharp corners, arc=0pt,
      left=6pt, right=6pt, top=2pt, bottom=2pt},
    overlay unbroken and first={
      \draw[KaitouBar, line width=\BarW] (frame.north west) -- (frame.south west);
    },
    arc=\ArcR, outer arc=0pt,
    left=12pt, right=10pt, top=14pt, bottom=8pt,
    fontupper=\color{Ink}\small,
    before skip=4pt, after skip=10pt,
  ]
}{%
  \end{tcolorbox}
}

% =============================================================
% 4択問題 (itpmon): ア/イ/ウ/エ 本試験準拠
% =============================================================
\newcounter{itpmonctr}
\renewcommand{\theitpmonctr}{\arabic{itpmonctr}}

% #1 = key-value options, #2 = problem text
\NewDocumentEnvironment{itpmon}{O{} m}{%
  \refstepcounter{itpmonctr}%
  \par\addvspace{8pt}%
  \begin{tcolorbox}[
    enhanced, breakable,
    colback=white, colframe=HairLine,
    boxrule=0.5pt, arc=0pt, outer arc=0pt,
    left=10pt, right=10pt, top=10pt, bottom=8pt,
    before skip=8pt, after skip=4pt,
  ]
    \noindent
    \tikz[baseline=(q.base)]{
      \node[fill=QuizBar, text=white, inner xsep=7pt, inner ysep=2pt,
            font=\sffamily\scriptsize\bfseries, rounded corners=1pt] (q)
        {Q\arabic{itpmonctr}};
    }%
    \hspace{4pt}%
    {\small\color{Muted}#1}%
    \par\vspace{5pt}%
    \textbf{\color{Navy}#2}\par
    \vspace{4pt}%
}{%
  \end{tcolorbox}
}

% 選択肢マクロ: 4つの選択肢を表示
\NewDocumentCommand{\sentakushi}{m m m m}{%
  \begin{enumerate}[label=\textbf{\color{SubBlue}\ifcase\value{enumi}\or ア\or イ\or ウ\or エ\fi.},
                   leftmargin=2.2em, itemsep=3pt, topsep=4pt]
    \item #1
    \item #2
    \item #3
    \item #4
  \end{enumerate}
}

% 解説マクロ: 正解＋解説
% #1 = 正解 (ア/イ/ウ/エ)
\NewDocumentEnvironment{kaisetsu}{m}{%
  \begin{kaisetsubox}
    \textbf{\color{KaiBar}正解}\quad\textbf{#1}\par
    \vspace{3pt}%
}{%
  \end{kaisetsubox}
}

% =============================================================
% 到達確認チェックリスト
% =============================================================
\NewDocumentEnvironment{tatsaku}{O{到達確認チェックリスト}}{%
  \par\addvspace{12pt}%
  \begin{tcolorbox}[
    enhanced, breakable,
    colback=white, colframe=Teal,
    boxrule=0.6pt, arc=2pt, leftrule=4pt,
    left=12pt, right=10pt, top=10pt, bottom=10pt,
    before skip=12pt, after skip=12pt,
    title={\sffamily\small\bfseries #1},
    colbacktitle=Teal, coltitle=white,
    attach boxed title to top left={xshift=-4pt, yshift=-6pt},
    boxed title style={colframe=Teal, boxrule=0pt, sharp corners, arc=0pt,
      left=8pt, right=8pt, top=2pt, bottom=2pt},
  ]
  \small\color{Ink}%
  \begin{itemize}[leftmargin=1.8em, itemsep=3pt, topsep=2pt,
    label={\color{Teal}$\square$}]
}{%
  \end{itemize}
  \end{tcolorbox}
}

% =============================================================
% タイムアタック
% =============================================================
\NewDocumentEnvironment{timeattack}{m m m}{%
  \par\addvspace{14pt}%
  \noindent
  \tikz[baseline=(t.base)]{\node[fill=Navy, text=white, inner xsep=8pt, inner ysep=3pt,
    font=\sffamily\scriptsize\bfseries, rounded corners=1pt] (t) {TIME ATTACK};}%
  \hspace{4pt}%
  \tikz[baseline=(u.base)]{\node[fill=Teal, text=white, inner xsep=6pt, inner ysep=2pt,
    font=\sffamily\scriptsize\bfseries, rounded corners=1pt] (u) {#2};}%
  \hspace{4pt}%
  \tikz[baseline=(v.base)]{\node[fill=BadgeHyojun, text=white, inner xsep=6pt, inner ysep=2pt,
    font=\sffamily\scriptsize\bfseries, rounded corners=1pt] (v) {目標 #3};}%
  \par\vspace{6pt}%
  {\Large\bfseries\color{Navy}#1\par}%
  \vspace{3pt}%
  {\color{HairLine}\rule{\linewidth}{0.4pt}}\par
  \vspace{6pt}%
}{%
  \par\vspace{6pt}%
  {\color{HairLine}\rule{\linewidth}{0.4pt}}\par
  \vspace{10pt}%
}

% =============================================================
% 模試形式セット
% =============================================================
\NewDocumentEnvironment{mogishi}{m m m}{%
  \clearpage
  \noindent
  \tikz[baseline=(m.base)]{\node[fill=Navy, text=white, inner xsep=10pt, inner ysep=4pt,
    font=\sffamily\small\bfseries, rounded corners=1pt] (m) {模試形式};}%
  \hspace{4pt}%
  \tikz[baseline=(n.base)]{\node[fill=Teal, text=white, inner xsep=8pt, inner ysep=3pt,
    font=\sffamily\scriptsize\bfseries, rounded corners=1pt] (n) {#2};}%
  \hspace{4pt}%
  \tikz[baseline=(o.base)]{\node[fill=BadgeHyojun, text=white, inner xsep=8pt, inner ysep=3pt,
    font=\sffamily\scriptsize\bfseries, rounded corners=1pt] (o) {目標 #3};}%
  \par\vspace{8pt}%
  {\color{TopBar}\rule{0.4\linewidth}{4pt}}%
  {\color{AccentBar}\rule{0.2\linewidth}{4pt}}\par
  \vspace{8pt}%
  {\LARGE\bfseries\color{Navy}#1\par}%
  \vspace{6pt}%
  {\color{HairLine}\rule{\linewidth}{0.4pt}}\par
  \vspace{8pt}%
}{%
  \par\vspace{10pt}%
  {\color{TopBar}\rule{\linewidth}{1pt}}\par
  \vspace{12pt}%
}

% =============================================================
% まとめボックス
% =============================================================
\MakeActorBox{matome}{この章のまとめ}{Navy}{LeadBG}{Navy}
